\section{Threats to Validity}
\label{sec:threats}

\emph{Construct validity.} The central construct is mission-critical exposure, which we
operationalize through the Mission-Critical True Positive target and the Mission-Weighted
Precision metric. Because both the target and CAP-G reference the Asset Context Score, a reader
must ask whether the advantage is a measurement of mission precision or an artifact of a shared
term. We mitigate this circularity three ways. The MCTP target is fixed to the primary weighted
ACS and applied identically to every method, including the ablations and the baselines, so no
ranker gains a labelling advantage. The homogeneous-fleet control (H3) shows that a shared term
cannot by itself manufacture an advantage, since the realized delta there is exactly zero. And
the context-blind tradeoff check (H4) confirms that CAP-G does not dominate the
\emph{independent} hygiene-defined target the base scorer was tuned for, so the gain on the
mission target is not bought by degrading the original one. The pre-registered
circularity-investigation rule, which would have triggered had CAP-G dominated both opposing
targets, was checked and not triggered. A residual construct risk remains in the definition of
the MCTP threshold itself ($\mathrm{EPSS}>0.10$ and ACS above the fleet 75th percentile), which
is a reasonable but not unique encoding of what mission-critical means.

\emph{Internal validity.} The study is pre-registered: hypotheses, thresholds, metrics, and the
evaluation seed range were fixed in a dated protocol before any evaluation-seed result was
inspected~\cite{wohlin2012experimentation}, and calibration used five disjoint held-out seeds, so
no analytic choice could be steered by the outcome. The calibrated context-emphasis $\rho$
saturated the locked grid; per the pre-registered stop rule we did not expand the grid post hoc,
so the evaluation uses a conservative lower bound on context emphasis and the reported gains are
if anything understated with respect to $\rho$, not cherry-picked. Every reported number is read
mechanically from the frozen results files rather than hand-entered, which removes a transcription
path for error.

\emph{External validity.} Hosts, context attributes, and fleet structure are synthetic with
public priors drawn from FIPS~199, the CISA directives, and the CJIS policy; only the EPSS and KEV
CVE signal is real, resampled from the vendored 203{,}174-CVE corpus and identical in distribution
across all compared methods. The category shares are documented priors, not a census of any
agency, so the absolute MWP values should be read as illustrative of the model rather than as
field estimates; the comparative deltas are the result. A single fleet scale ($\sim$830 hosts) was
evaluated because capacity, not host count, is the swept axis, and the heterogeneity mechanism
does not depend on the specific scale. Different context mixes would shift magnitudes but not the
qualitative findings, which follow from the structure of the model rather than its constants;
calibrating the mixes against a real agency asset inventory is the natural next step and would
convert the illustrative magnitudes into estimates.

\emph{Statistical validity.} All intervals are 95\% bias-corrected and accelerated bootstrap
intervals over 25 seeds with 10{,}000 resamples~\cite{efron1987bca,efron1994introduction}, a
choice appropriate to the small-sample, possibly skewed statistics reported here. The evidentiary
standard throughout is interval non-overlap rather than null-hypothesis significance testing,
which avoids the multiple-comparison pitfalls of repeated $p$-values across many hypotheses,
capacities, and ablations. The primary hypothesis is reported as partial precisely because the
realized capacity-averaged advantage of $4.4$~pp falls below the locked $5$~pp bar even though
every individual interval excludes zero, which is the conservative reading the interval-based
standard requires.
